\documentclass[10pt,a4paper]{amsart}
\usepackage[margin=19mm]{geometry}
\usepackage{amsmath,amssymb,mathtools}
\usepackage[scr=rsfs,cal=euler]{mathalfa}
\usepackage{xfrac}
\usepackage[unicode,hidelinks,bookmarks=false]{hyperref}
\newcommand{\ie}{i.e.\ }
\newcommand{\cf}{cf.\ }
\newcommand{\resp}{respectively\ }
\newcommand{\Subsection}{\subsection}
\providecommand{\nobreakdash}{\nobreak\hbox{-}}
\setlength{\emergencystretch}{2em}
\multlinegap0pt
\usepackage{enumerate}
\usepackage{xcolor}
\allowdisplaybreaks
\newcommand{\meisam}[1]{{\textbf{\color{red}meisam: #1}}}

\newtheorem{Theorem}{Theorem}[section]
\newtheorem{Definition}[Theorem]{Definition}
\newtheorem{Proposition}[Theorem]{Proposition}
\newtheorem{Assumption}[Theorem]{Assumption}
\newtheorem{Lemma}[Theorem]{Lemma}
\newtheorem{Corollary}[Theorem]{Corollary}
\newtheorem{Remark}[Theorem]{Remark}
\newtheorem{Example}[Theorem]{Example}

\newcommand{\R}{\mathbb{R}}
\def \dd{{\rm d}}

\newcommand{\curly}[1]{\left\lbrace #1 \right\rbrace}
\newcommand{\bracket}[1]{\left[ #1 \right]}
\newcommand{\parenthesis}[1]{\left( #1 \right)}
\renewcommand{\vec}[1]{\mathbf{#1}}
\renewcommand{\O}{\Omega}
\renewcommand{\bf}[1]{\textbf{#1}}
\newcommand{\mcal}[1]{\mathcal{#1}}
\newcommand{\mbf}[1]{\mathbf{#1}}
\newcommand{\abs}[1]{\left |#1\right |}
\newcommand{\norm}[1]{\left \|#1\right \|}
\newcommand{\inpro}[1]{\left\langle#1\right\rangle}
\newcommand{\union}{\bigcup}
\newcommand{\intersection}{\bigcap}
\newcommand{\op}[1]{\operatorname{#1}}
\newcommand{\minus}{\setminus}
\newcommand{\nul}{\op{null}}
\newcommand{\range}{\op{range}}
\newcommand{\spn}{\op{span}}
\newcommand{\E}{\mathbb{E}}
\renewcommand{\P}{\mathbb{P}}

\setcounter{section}{1}
\begin{document}
\title[Regularity and well-posedness for the KL-regularized control problem]{Stochastic Control for Fine-tuning Diffusion Models: Optimality, Regularity, and Convergence}
\author{Yinbin Han, Meisam Razaviyayn, Renyuan Xu}
\date{Source: arXiv:2412.18164v4 (CC BY 4.0); Forty-Second International Conference on Machine Learning 2025}
\maketitle
\noindent\emph{Source and licence.} Stochastic Control for Fine-tuning Diffusion Models: Optimality, Regularity, and Convergence. Version arXiv:2412.18164v4 (CC BY 4.0); journal reference Forty-Second International Conference on Machine Learning 2025. This material is distributed under the Creative Commons Attribution 4.0 International licence, http://creativecommons.org/licenses/by/4.0/, as recorded in the arXiv OAI licence metadata for this exact version and shown on the abstract page https://arxiv.org/abs/2412.18164v4. Full rights notice: licenses/arxiv-2412.18164-CC-BY-4.0.txt.\par\smallskip
\noindent\textit{Editorial scope.} Counted unit: original Theorem 2.8 of the article, printed with its parenthetical title ``Regularity and well-posedness'' (native environment \texttt{Theorem}, source label \texttt{thm:well-posed}, source0.tex lines 276--289), the main regularity and well-posedness theorem announced in the abstract, delivered together with its complete original author proof as the single primary proof body at source lines 295--456. The proof is a backward induction over the denoising steps whose heading names its own statement (it opens with ``We prove Theorem \texttt{\textbackslash ref\{thm:well-posed\}} by backward induction''), so no detached proof block is involved. No proof marker is added and no mathematics is written, repaired or re-derived. Necessary complete context is retained uncounted: source lines 114--275, that is the whole of Section 2 ``Problem set-up and regularity results'' from its own section heading, with the two subsections ``Problem set-up'' and ``Regularity and well-posedness'', the recursive definition of the constants of the theory, Remark 2.1 to Remark 2.4, Lemma 2.5 with its complete proof, and Assumptions 2.6 and 2.7; and source lines 290--294, the author's discussion of the counted theorem and the brief proof sketch that the author prints immediately before the detailed proof. Nothing inside the delivered ranges is omitted, so every printed number of the delivered unit coincides with the version PDF: the delivered text is the whole of Section 2 and the wrapper restores the source's own counters editorially, so the section prints as Section 2, its subsections as 2.1 and 2.2, the theorem-like environments print as Remark 2.1, Remark 2.2, Remark 2.3, Remark 2.4, Lemma 2.5, Assumption 2.6, Assumption 2.7 and Theorem 2.8 exactly as in the source, and the numbered equations print with the source's own numbers, the first of them being the objective and dynamics equations (2) and (3). The publisher's own class is not present in the runtime; the article is delivered in the AMS amsart wrapper of the pinned TeX Live runtime together with the author's own macro preamble, the author's own theorem declarations, the author's own front matter and the macros of the author's own local style file \texttt{pkg.sty}, which the e-print keeps outside the main .tex and which the structure inventory records as unspliced; those macros are carried into the wrapper verbatim. The delivered unit begins at Section 2 rather than at the Introduction, because the whole of the Introduction is a literature review whose citations use the author's author-year \texttt{\textbackslash citet} command, which the corpus bibliography mechanism does not carry; every definition, assumption and equation that the counted statement and its proof use is defined inside the delivered Section 2, and the two forward references of the Introduction to later sections (the algorithm and discussion sections) are therefore not part of the unit. The article's own bibliography is not delivered as source text: the 14 works that the delivered ranges cite are re-typeset from the author's own \texttt{.bbl} in the author's own order under the published numbers of that bibliography, so the printed citation numbers coincide with the version PDF, and the BibTeX line-breaking command \texttt{\textbackslash newblock} of those entries, which only marks a line break inside a reference, is not reproduced. Printed slips kept literal, in particular the misspelling ``divergencd'' at source line 196, the misspelling ``Lipchitz'' at source line 293, the author's abbreviation $\E_{t+1\vert t}$ at source line 226 where the notation introduced two lines earlier is $\E_{p_{t+1\vert t}}$, and the middle dot of ``DALL·E'' at source line 134. No figure, table, drawing, algorithm float or external asset occurs anywhere in the delivered ranges, so no artwork of any kind is redistributed and no bounded source comparison was needed.
\section{Problem set-up and regularity results}
\label{sec:set-up}


\subsection{Problem set-up}
To address the challenge of fine-tuning, we first introduce the dynamics of denoising diffusion probabilistic models (DDPMs), a widely adopted pre-trained framework in practice. Our focus is on the discrete-time formulation, which aligns with the standard implementation of diffusion models in practice. In addition, DDPMs serve as a foundational scheme for effective fine-tuning.

\paragraph{Denoising diffusion probabilistic models (DDPM).} A well-trained DDPM $\{Y_t^{\rm pre}\}_{t=0}^T$ in discrete time usually follows the following stochastic dynamics with state $Y_t^{\rm pre}\in \mathbb{R}^d$:
\begin{align}
    Y_{t+1}^{\rm pre} & = \frac{1}{\sqrt{\alpha_t}}\big(Y_t^{\rm pre} + (1 - \alpha_t)s^{\rm pre}_t(Y_t^{\rm pre})\big) + \sigma_t W_t,   \;\; 0 \leq t \leq T - 1, \;\; \textit{with}\;\; Y_0^{\rm pre} \sim \mcal{N}(0, I_d),\label{eq:ddpm}
\end{align}
where $\curly{W_t}_{t = 0}^{T - 1}$ are i.i.d.~standard Gaussian random vectors such that $W_t\sim \mathcal{N}(0,I_d)$. The hyper-parameters $\curly{\alpha_t}_{t=0}^{T-1}$ with  $\alpha_t\in (0, 1)$ and $\curly{\sigma_t}_{t=0}^{T-1}$ with $\sigma_t>0$ represent the prescribed denoising rate schedules \cite{Ho20,li2023towards}. They control the variance of noise in data generation\footnote{While DDPM chooses $\sigma_t^2 = 1/\alpha_t - 1$, our analysis in the subsequent sections works for general $\sigma_t$.}.  Here, $s^{\rm pre}_t:\mathbb{R}^d \rightarrow\mathbb{R}^d$ is the score function associated with the pre-trained model. In practice, the score estimator $s_t^{\rm pre}$  is obtained by training a neural network to minimize the score matching loss \cite{hyvarinen2005estimation}. {With a well-trained pre-trained model, we expect the distribution of $Y_T$ to be close to the target distribution, from which the pre-trained model has access to samples and seeks to generate additional samples.} Given $s_t^{\rm pre}$, we denote $ p_{t+1 \vert t}^{\rm pre}(\cdot \vert y_t) $ as the conditional density of $Y_{t+1}^{\rm pre}$ given $Y_t^{\rm pre} = y_t$ induced by the dynamics \eqref{eq:ddpm}, i.e., 
\begin{align*}
    p_{t+1\vert t}^{\rm pre}({}\cdot{}\vert y_{t}) = f\parenthesis{{}\cdot{}\bigg\vert\frac{1}{\sqrt{\alpha_t}}\parenthesis{y_t + (1 - \alpha_t)s^{\rm pre}_t(y_t)}, \sigma_{t}^{2} I_d}.
\end{align*}


%

\begin{Remark}[Choice of the pre-trained model.]
    DDPMs have become the leading choice for pre-trained models across a wide range of applications, serving as a powerful building block for fine-tuning in diverse tasks \cite{dhariwal2021diffusion,Ho20,watson2023novo,li2023towards}.  By pre-training on large datasets, DDPMs capture complex data distributions, enabling relatively easier fine-tuning for specialized applications such as conditional image synthesis, text-to-image generation, and even time-series forecasting. This versatility has been demonstrated in models such as Stable Diffusion \cite{rombach2022high} and DALL·E \cite{ramesh2022hierarchical}, where fine-tuning on task-specific data enhances performance and customization \cite{fan2024reinforcement}. 
\end{Remark}

\paragraph{Stochastic control formulation.}
When human preference or soft constraint can be modeled by a (stochastic) reward function $R(\cdot)$, we consider the following  stochastic control problem to model fine-tuning tasks of diffusion models:
\begin{eqnarray}
    \sup_{\curly{U_t}_{t = 0}^{T-1}\in \mathcal{U}_0} &  \mathbb{E}\left[R(Y_T) - \sum_{t = 0}^{T-1}\beta_t{\rm KL}\parenthesis{p_{t+1\vert t} ({}\cdot{}|Y_t)\,\Big\|\,p_{t+1\vert t}^{\rm pre}({}\cdot{}|Y_t)}\right], \label{eq:objective-dst} \\ 
    \text{s.t.} & \;\; 
    Y_{t+1} = \frac{1}{\sqrt{\alpha_t}}\parenthesis{Y_t + (1 - \alpha_t)U_t} + \sigma_t W_t,\quad y_{0} \sim \mathcal{N}(0, I_{d}), \label{eq:dynamics-dst}
\end{eqnarray}
   with state dynamics $Y_t\in \mathbb{R}^d$ and control $U_t\in \mathbb{R}^d$. In particular, we work with Markovian policies such that the objective in \eqref{eq:objective-dst} is well-defined. Mathematically, we specify the admissible control set as 
   \begin{align*}
    \mathcal{U}_{t}:=\Big\{\{U_\ell\}_{\ell=t}^{T-1} \,\Big|& U_\ell = {u_\ell(Y_\ell)} \text{ for some measurable function } u_\ell: \mathbb{R}^d\rightarrow \mathbb{R}^d \\
    & \text{ and }\mathbb{E}\Big[R(Y_T) - \sum_{\ell = t}^{T-1}\beta_t{\rm KL}\Big(p_{\ell+1\vert \ell} ({}\cdot{}|Y_\ell)\,\Big\|\,p_{\ell+1\vert \ell}^{\rm pre}({}\cdot{}|Y_\ell)\Big)\Big]<\infty\Big\}, \; 0 \leq t \leq T-1.
   \end{align*}
In terms of the reward, we assume $R\in L^1(\mathbb{R}^d,\mathcal{F}_0)$, namely, $R(y)\in \mathcal{F}_0$ has finite mean for all $y\in \mathbb{R}^d$.
%
   We assume the reward function is {\it known} in this work.
   The KL-divergence measures the deviation of the fine-tuned model $p_{t+1\vert t}$ from the pre-trained model $p_{t+1\vert t}^{\rm pre}$, where $ p_{t+1 \vert t} $ denotes the conditional distribution of $ Y_{t+1} $ given $ Y_{t} $ induced by the control policy $u_t$. The regularization coefficients $\{\beta_t\}_{t=0}^{T-1}$ control the strength of the regularization term.  
   %


{We have a few remarks in place.}    
\begin{Remark}[Control as score function]
    Our formulation is rooted in both diffusion models and control theory. Comparing \eqref{eq:dynamics-dst} with \eqref{eq:ddpm}, we replace the pre-trained score $s_t^{\rm pre}$ by a control variable $u_t$. Consequently, in the context of diffusion models, the learned control sequence $\{u_t\}_{t=0}^{T-1}$ can be viewed as the {\it new} score function of the fine-tuned model. In other words, solving the control problem \eqref{eq:objective-dst}-\eqref{eq:dynamics-dst} is essentially learning a new score function in response to the reward signal $R(\cdot)$ for fine-tuning. Note that the linear dynamics \eqref{eq:dynamics-dst} with Gaussian noise is preferred in stochastic control due to its tractability in analysis. Specifically, our theoretical analysis in the subsequent sections heavily relies on the linearity and the Gaussian smoothing~effect.
\end{Remark} 

\begin{Remark}[Rationale for the objective function and the known reward assumption]
     The objective function in the control formulation\eqref{eq:objective-dst}-\eqref{eq:dynamics-dst}  consists of two parts:~a terminal reward function $R(\cdot)$ at time $T$ and intermediate KL penalties. The reward $R(\cdot)$ captures human preference on the generated samples. For example, in the text-to-image generation, the reward $R$ represents how the generated data $Y_T$ is aligned with the input prompt \cite{fan2024reinforcement,black2023training}. In addition, the penalty term ensures that the fine-tuned model is not too far away from the pre-trained model, which prevents overfitting. Unlike the Shannon entropy of the (randomized) control policy commonly used in the RL literature to encourage exploration, the KL regularization in \eqref{eq:objective-dst}  is applied between two conditional probability densities. From an optimization perspective, the KL divergence introduces   concavity/convexity in control and consequently leads to a better landscape of the objective. Indeed, we will choose the parameter $\beta_t$ sufficiently large to guarantee the existence and uniqueness of the optimal control and to satisfy certain regularity conditions; see Theorem \ref{thm:well-posed}.

    In practice, fine-tuning is typically performed on a dataset containing reward ratings for each data sample, which is much smaller than the pre-trained dataset. Using this fine-tuning dataset, one can estimate the expected reward function. The assumption of known reward is not overly restrictive, as online learning to acquire new rewards is generally uncommon.
\end{Remark}    
    
   %
    
  
    \begin{Remark}[Connection to KL divergence on path-wise measures.]
    Another natural idea is to impose the KL divergence between the path-wise measures, i.e., $ {\rm KL}(p_{0:T} \| p_{0:T}^{\rm pre}) $, where $ p_{0:T} $ is the joint density of $ \curly{Y_{t}}_{t = 0}^{T} $ and $ p_{0:T}^{\rm pre} $ is the joint density of $ \curly{Y_{t}^{\rm pre}}_{t = 0}^{T} $. This choice is relevant to formulation \eqref{eq:dynamics-dst}. In particular, the Markov property of the dynamics and the chain rule of the KL divergence imply
    \begin{align*}
        {\rm KL}(p_{0:T} \| p_{0:T}^{\rm pre}) = \E\bracket{\sum_{t = 0}^{T-1}{\rm KL}\Big(p_{t+1\vert t}({}\cdot{}\vert Y_t) \|p_{t+1\vert t}^{\rm pre}({}\cdot{}\vert Y_t)\Big)},
    \end{align*}
    where the expectation is taken over all random variables $\{Y_t\}_{t = 0}^T$. Here, we define
    \begin{align}
        {\rm KL}(p_{0:T} \| p_{0:T}^{\rm pre}) & \coloneqq \int p_{0:T}(y_0, \dots, y_t)\log\frac{p_{0:T}(y_0, \dots, y_t)}{p_{0:T}^{\rm pre}(y_0, \dots, y_t)}{\rm d}y_0\cdots {\rm d}y_t, \label{eq:kl-path}
    \end{align} 
    and for given $y_t \in \R^d$, denote
    \begin{align}
        {\rm KL}\Big(p_{t+1\vert t}({}\cdot{}\vert y_t) \|p_{t+1\vert t}^{\rm pre}({}\cdot{}\vert y_t)\Big) & \coloneqq \int p_{t+1\vert t}(y_{t+1} \vert y_t) \log \frac{p_{t+1\vert t}(y_{t+1} \vert y_t)}{p^{\rm pre}_{t+1\vert t}(y_{t+1} \vert y_t)}{\rm d}y_{t+1}.  \label{eq:kl-condition}
    \end{align}
%
Thus, when $\beta_t = \beta$ for all $t$, the objective in \eqref{eq:dynamics-dst} becomes
\begin{align*}
    \mathbb{E}\left[R(Y_T) - \sum_{t = 0}^{T-1}\beta_t{\rm KL}\Big(p_{t+1\vert t}({}\cdot{}\vert Y_t) \|p_{t+1\vert t}^{\rm pre}({}\cdot{}\vert Y_t)\Big)\right] = \mathbb{E}\left[R(Y_T)\right] - \beta  {\rm KL}(p_{0:T} \| p_{0:T}^{\rm pre}),
\end{align*}
which is a common choice in fine-tuning of diffusion models \cite{fan2024reinforcement,black2023training,zhang2001some,gao2024reward}. 
%
    \end{Remark}
   
  Given two Gaussian densities $p_{t+1\vert t}$ and $p^{\rm pre}_{t+1\vert t}$ in  \eqref{eq:dynamics-dst}, the following lemma simplifies the KL divergence term.

%
  \begin{Lemma} \label{lemma:kl-l2} 
  Consider the KL divergencd defined as in \eqref{eq:kl-condition}.
  For any $y_t \in \R^d$ and any admissible control policy $u_t$, it holds that
      \begin{align}\label{eq:kl-l2-1}
        {\rm KL}\Big(p_{t+1\vert t}({}\cdot{}\vert y_t) \|p_{t+1\vert t}^{\rm pre}({}\cdot{}\vert y_t)\Big) = \frac{(1 - \alpha_{t})^{2}}{2\alpha_{t}\sigma_{t}^{2}}\norm{u_{t}(y_t) - s^{\rm pre}_{t}(y_{t})}_{2}^{2}.
    \end{align}
  \end{Lemma}
    
  Lemma \ref{lemma:kl-l2} links the KL-divergence between two conditional distributions with the squared loss between the control $u_t$ and the pre-trained score $s_t^{\rm pre}$. As the control can be interpreted as the new score of the fine-tuned model, \eqref{eq:kl-l2-1} enjoys the spirit of the score matching loss in training diffusion models \cite{Ho20,song2020score,han2024neural}. The proof of Lemma \ref{lemma:kl-l2} is based on direct calculations.

    \begin{proof}[Proof of Lemma \ref{lemma:kl-l2}]
        Recall that for any $y_t \in \R^d$,
        \begin{align*}
           {p_{t+1\vert t}({}\cdot{}\vert y_{t})} = f\parenthesis{{}\cdot{}\vert\mu_t(y_t), \sigma_{t}^{2} I_d} \;\; \text{and}\;\; p_{t+1\vert t}^{\rm pre}({}\cdot{}\vert y_{t}) = f\parenthesis{{}\cdot{}\vert\mu_t^{\rm pre}(y_t), \sigma_{t}^{2} I_d},
        \end{align*}
        with 
        \begin{align*}
            \mu_t(y_t) = \frac{1}{\sqrt{\alpha_t}}\parenthesis{y_t + (1 - \alpha_t)u_t(y_t)} \;\; \text{and} \;\; \mu_t^{\rm pre}(y_t) = \frac{1}{\sqrt{\alpha_t}}\parenthesis{y_t + (1 - \alpha_t)s^{\rm pre}_t(y_t)}.
        \end{align*}
        Thus, for any $y_t, y_{t+1} \in \R^d$, we have
        \begin{align*}
            \log \left(\frac{p_{t+1\vert t}(y_{t+1} \vert y_t)}{p^{\rm pre}_{t+1\vert t}(y_{t+1} \vert y_t)} \right)= -\frac{1}{2\sigma_t^2}\norm{y_{t+1} - \mu_t(y_t)}_2^2 + \frac{1}{2\sigma_t^2}\norm{y_{t+1} - \mu_t^{\rm pre}(y_t)}_2^2. 
        \end{align*}
        Denote $\E_{p_{t+1\vert t}}$ as the expectation under the conditional density $p_{t+1\vert t}(\cdot\vert y_t)$ of $Y_{t+1}$ given $Y_t = y_t$. By definition of the KL divergence, we have %
        \begin{align}
             {\rm KL}\Big(p_{t+1\vert t}({}\cdot{}\vert y_t) \|p_{t+1\vert t}^{\rm pre}({}\cdot{}\vert y_t)\Big) & =\E_{p_{t+1\vert t}}\bracket{\log \left(\frac{p_{t+1\vert t}(Y_{t+1} \vert y_t)}{p^{\rm pre}_{t+1\vert t}(Y_{t+1} \vert y_t)}\right)} \nonumber\\ & = -\frac{1}{2\sigma_t^2}\E_{p_{t+1\vert t}}\bracket{\norm{Y_{t+1} - \mu_t(y_t)}_2^2} +  \frac{1}{2\sigma_t^2}\E_{p_{t+1\vert t}}\bracket{\norm{Y_{t+1} - \mu_t^{\rm pre}(y_t)}_2^2}. \label{eq:kl-simple}
        \end{align}
        Note that 
        \begin{align*}
            \E_{p_{t+1\vert t}}\bracket{\norm{Y_{t+1} - \mu_t^{\rm pre}(y_t)}_2^2} & = \E_{p_{t+1\vert t}}\bracket{\norm{Y_{t+1} - \mu_t(y_t)}_2^2} + \E_{p_{t+1\vert t}}\bracket{\norm{\mu_t(y_t) - \mu_t^{\rm pre}(y_t)}_2^2} \\ & \quad + 2\E_{p_{t+1\vert t}}\bracket{(Y_{t+1} - \mu_t(y_t))^\top(\mu_t(y_t) -\mu_t^{\rm pre}(y_t))} \\ & = \E_{p_{t+1\vert t}}\bracket{\norm{Y_{t+1} - \mu_t(y_t)}_2^2} + \norm{\mu_t(y_t) - \mu_t^{\rm pre}(y_t)}_2^2,
        \end{align*}
        where we use the fact that $\E_{t+1\vert t}\bracket{Y_{t+1}} = \mu_t(y_t)$.
        Plugging the above equality into \eqref{eq:kl-simple}, we obtain
        \begin{align*}
            {\rm KL}\Big(p_{t+1\vert t}({}\cdot{}\vert y_t) \|p_{t+1\vert t}^{\rm pre}({}\cdot{}\vert y_t)\Big) = \frac{1}{2\sigma_t^2}\norm{\mu_t(y_t) - \mu_t^{\rm pre}(y_t)}_2^2 =  \frac{(1 - \alpha_{t})^{2}}{2\alpha_{t}\sigma_{t}^{2}}\norm{u_{t}(y_t) - s^{\rm pre}_{t}(y_{t})}_{2}^{2},
        \end{align*}
        which completes the proof. 
    \end{proof}

    
    With Lemma \ref{lemma:kl-l2}, we define the optimal value function at time $ t $ as
    \begin{align}
    V^{*}_{t}(y) &\coloneqq \sup_{\curly{U_{\ell}}_{\ell \geq t}\in \mcal{U}_t}\mathbb{E}\left[R(Y_T) - \sum_{\ell = t}^{T-1}\beta_\ell{\rm KL}\Big(p_{\ell+1\vert \ell}({}\cdot{}\vert Y_\ell) \|p_{\ell+1\vert \ell}^{\rm pre}({}\cdot{}\vert Y_\ell)\Big)\,\Big\vert \,Y_t=y\right] \nonumber \\
    &=\sup_{\curly{u_{\ell}}_{\ell \geq t}}\mathbb{E}\left[R(Y_T) - \sum_{\ell = t}^{T-1}\beta_\ell\frac{(1 - \alpha_{\ell})^{2}}{2\alpha_{\ell}\sigma_{\ell}^{2}}\norm{u_{\ell}(Y_\ell) - s^{\rm pre}_{\ell}(Y_{\ell})}_{2}^{2}\,\Big\vert \,Y_t=y\right]. \label{eq:opt-value}
    \end{align}
The Dynamic Programming Principle implies that $V_t^*$ satisfies the Bellman equation:
    \begin{align}
        V^{*}_{t}(y) =\sup_{u_{t}} \ \E\bracket{V^{*}_{t+1}\left(\frac{1}{\sqrt{\alpha_t}}\parenthesis{y + (1 - \alpha_t)u_t(y)} + \sigma_t W_t\right) - \beta_t \frac{(1 - \alpha_{t})^{2}}{2\alpha_{t}\sigma_{t}^{2}}\norm{u_{t}(y) - s^{\rm pre}_{t}(y)}_{2}^{2}}.\label{eq:bellman-V*}
    \end{align} 
In the next section, we will discuss the well-posedness of the optimal control problem \eqref{eq:objective-dst}--\eqref{eq:dynamics-dst} and the regularity of the optimal value function $V_t^*$. 
    
\subsection{Regularity and well-posedness}
In this section, we establish the regularity of the optimal value function and the optimal control policy. To start, we assume the following assumptions on the reward and pre-trained score functions.
\begin{Assumption}[Smoothness of the reward]\label{ass:smooth r}
    Assume $r(y) \coloneqq \mathbb{E}[R(y)]$ is $L_{0}^r$-Lipschitz and $L_{1}^r$-gradient Lipschitz in $y \in \R^d$, i.e.,
    \begin{align}
        \abs{r(y_1) - r(y_2) } \leq L_{0}^r\norm{y_1 - y_2}_2, \\ 
        \norm{\nabla r(y_1) - \nabla r(y_2) }_2 \leq L_{1}^r\norm{y_1 - y_2}_2.
    \end{align}
\end{Assumption}

\begin{Assumption}[Smoothness of the pre-trained score function]\label{ass:smooth s}
    Assume $s_t^{\rm pre}$ is $L_{0}^{s}$-Lipschitz and $L_{1}^{s}$-gradient Lipschitz in $y \in \R^d$, i.e.,
    \begin{align}
        \norm{s_t^{\rm pre}(y_1) - s_t^{\rm pre}(y_2) }_2 \leq L_{0,t}^{s}\norm{y_1 - y_2}_2, \\ 
        \norm{\nabla s_t^{\rm pre}(y_1) - \nabla s_t^{\rm pre}(y_2) }_2 \leq L_{1,t}^{s}\norm{y_1 - y_2}_2.
    \end{align}
\end{Assumption}

While Assumptions \ref{ass:smooth r} and \ref{ass:smooth s} guarantee the smoothness of both the expected reward $r$ and the pre-trained score $\curly{s_t^{\rm pre}}_{t=0}^{T-1}$, no convexity assumptions are imposed, and thus the control problem \eqref{eq:objective-dst}--\eqref{eq:dynamics-dst} is in general {\it non-concave} with respect to $\curly{u_t}_{t=0}^{T-1}$.

The next theorem establishes the existence and uniqueness of the optimal control, as well as the regularities of the solution to problem \eqref{eq:objective-dst}--\eqref{eq:dynamics-dst}. For ease of exposition, we define  a series of constants recursively. Let $0 < \lambda_t < 1$ for each $ t < T$. Set $L_{0, T}^{V^{*}} = L_0^r$ and $L_{1, T}^{V^{*}} = L_1^r$. For $ t < T$, define
\begin{eqnarray}
    L_{0,t}^{V^{*}} &=& \frac{1}{\sqrt{\alpha_t}}(1 + (1 - \alpha_t)L_{0, t}^s)L_{0, t+1}^{V^{*}}, \label{eq:LV0star}\\ 
    L_{1, t}^{V^{*}} &=& \frac{1}{\alpha_{t}}(1 + (1 - \alpha_{t})L_{0, t}^s)(1 + (1 - \alpha_{t})L_{0, t}^{u^{*}})L_{1,t+1}^{V^{*}} + \frac{1 - \alpha_{t}}{\sqrt{\alpha_{t}}}L_{1, t}^s L_{0,t+1}^{V^{*}},\label{eq:LV1star} \\ 
    L_{0, t}^{u^*} & =& \lambda_{t}^{-1}\parenthesis{L_{0, t}^s + \frac{1 - \lambda_t}{1- \alpha_t}}, \label{eq:Lu0star}\\ 
    L_{1, t}^{u^*}  &= & \lambda_{t}^{-1}\parenthesis{L_{1, t}^{s} +  \frac{\E\bracket{\norm{W_{t}}_2}(1 - \lambda_t)}{(1 - \alpha_t)\sqrt{\alpha_{t}}\sigma_{t}}\parenthesis{1 + (1 - \alpha_{t})L_{0,t}^{u^{*}}}^{2}}.\label{eq:Lu1star}
\end{eqnarray}
Here, $\E\bracket{\norm{W_t}_2} = \frac{\sqrt{2}\Gamma((d+1)/2)}{\Gamma(d/2)} < \infty $ is a constant with $\Gamma(\cdot)$ being the Gamma function. Note that the above constants only depend on the system parameters $\alpha_t$, $L_0^r$, $L_1^r$, $L_0^s$ and $L_1^s$ as well as the hyper-parameter $\lambda_t$. In the next theorem, we show that the above constants are indeed the Lipschitz coefficients for $\curly{V_t^{*}}_{t = 0}^{T}$ and $\curly{u_t^{*}}_{t=0}^{T-1}$.
%


\begin{Theorem}[Regularity and well-posedness]\label{thm:well-posed}
Suppose Assumptions \ref{ass:smooth r} and \ref{ass:smooth s} hold and let the constants be defined according to \eqref{eq:LV0star}--\eqref{eq:Lu1star}. Choose $\beta_t$ such that $ 1 - \frac{\sigma_t^2}{\beta_t}L_{1, t+1}^{V^*} \geq \lambda_{t} > 0  $. Then,
\begin{enumerate}[(i)]
\item The optimal value function $V_t^*(y)$ defined in \eqref{eq:opt-value} is $L_{0,t}^{V^{*}}$-Lipschitz, differentiable and $L_{1,t}^{V^{*}}$-gradient Lipschitz in $y \in \R^d$.

    \item There is a unique optimal control $u_t^*: \R^d \to \R^d$ of problem \eqref{eq:objective-dst}--\eqref{eq:dynamics-dst} satisfying
\begin{eqnarray}\label{eq:u*(y)}
     u_{t}^{*}(y)  = s_{t}^{\rm pre}(y) + \frac{\sqrt{\alpha_{t}}\sigma_{t}^{2}}{(1 - \alpha_{t})\beta_t}\E\bracket{\nabla V_{t+1}^{*}\left(\frac{1}{\sqrt{\alpha_t}}\parenthesis{y + (1 - \alpha_t)u_t^*(y)} + \sigma_t W_t\right)}.
\end{eqnarray}
Moreover, the optimal value function $V_t^*$ is the unique $\mcal{C}^1$ solution to the Bellman equation \eqref{eq:bellman-V*}. 
\item \label{item:lip-u*} The optimal control $u_t^*(y)$ is $L_{0, t}^{u^{*}}$-Lipschitz, differentiable, and $ L_{1,t}^{u^{*}} $-gradient Lipschitz in $y \in \R^d$.

\end{enumerate}
\end{Theorem}


When the regularization coefficient $\beta_t$ is sufficiently large, Theorem \ref{thm:well-posed} states that the optimal value function $V_t^*(\cdot)$ is Lipschitz and gradient Lipschitz. The choice of $\beta_t$ also ensures the right hand side of \eqref{eq:bellman-V*} is strongly concave in $u_t(\cdot)$, guaranteeing the existence and uniqueness of the optimal control $u_t^*$. Furthermore, Theorem \ref{thm:well-posed} shows the optimal control $u_t^*$ is also Lipschitz and gradient Lipschitz. The regularities of $V_t^*$ and $u_t^*$ will serve as the foundation of the algorithm design and convergence analysis in the subsequent sections. \\

{We outline a brief proof sketch to highlight the ideas before providing the detailed proof, which is based on backward induction. Assuming $V_{t+1}^*$ is Lipschitz and gradient Lipschitz, the first-order optimality condition  implies that \eqref{eq:u*(y)} holds, given the  optimization problem is unconstrained. Direct calculation leads to the Lipchitz condition of $u_t^*$ and the smoothing effect of the Gaussian random variable is utilized to obtain differentiability. In addition, the Lipschitz property of $\nabla u_t^*$ is a consequence of the integration by parts formula in \eqref{eq:magic-1}; hence \eqref{item:lip-u*} holds. Finally, the regularities of $V_t^*$ follow from the Lipschitz conditions of $u_t^*$ and $\nabla u_t^*$, which completes the induction.}


\begin{proof}
We prove Theorem \ref{thm:well-posed} by backward induction. At time $t = T$, Assumption \ref{ass:smooth r} implies $V_T^*(y) = R(y)$ is $L_{0, T}^{V^*}$-Lipschitz and $L_{1, T}^{V^*}$-gradient Lipschitz with $L_{0, T}^{V^*} = L_0^r$ and $L_{1, T}^{V^*} = L_1^r$.

\underline{Step 1: Lipschitz condition of $ u_{t}^{*} $.} Assume that $V^*_{t+1}$ is $L_{0, t+1}^{V^*}$-Lipschitz and $L_{1, t+1}^{V^*}$-gradient Lipschitz in~$y\in \R^d$. The choice of $\beta_t$ implies that the mapping
\begin{eqnarray}\label{eq:obj-concave}
    u \mapsto \E\bracket{V^{*}_{t+1}\left(\frac{1}{\sqrt{\alpha_t}}\parenthesis{y + (1 - \alpha_t)u} + \sigma_t W_t\right) - \beta_t \frac{(1 - \alpha_{t})^{2}}{2\alpha_{t}\sigma_{t}^{2}}\norm{u - s^{\rm pre}_{t}(y)}_{2}^{2}}
\end{eqnarray}
is $\gamma_t$-strongly concave with $ \gamma_{t} = \frac{(1 - \alpha_{t})^{2}}{\alpha_{t}}\parenthesis{\frac{\beta_{t}}{\sigma_{t}^{2}} - L_{1, t+1}^{V^{*}}} > 0 $. Hence, there is a unique optimal control $u_t^*$ satisfying
\begin{eqnarray}\label{eq:u*(y)-re}
     u_{t}^{*}(y)  = s_{t}^{\rm pre}(y) + \frac{\sqrt{\alpha_{t}}\sigma_{t}^{2}}{(1 - \alpha_{t})\beta_t}\E\bracket{\nabla V_{t+1}^{*}\left(\frac{1}{\sqrt{\alpha_t}}\parenthesis{y + (1 - \alpha_t)u_t^*(y)} + \sigma_t W_t\right)},
\end{eqnarray}
which is obtained by setting the gradient of mapping \eqref{eq:obj-concave} as zero for each $y \in \R^d$.
Here, we apply the Lipschitz condition of $V_{t+1}^*$ and the dominated convergence theorem to interchange the operators. 

We next prove the Lipschitz condition of $u_t^*$ in $y\in \R^d$. Note that for any $y_1$ and $y_2$ in $\R^d$, Eq.~\eqref{eq:u*(y)-re} implies
\begin{align*}
    u_t^*(y_1) - u_t^*(y_2) & = s_{t}^{\rm pre}(y_1) - s_{t}^{\rm pre}(y_2)  + \frac{\sqrt{\alpha_{t}}\sigma_{t}^{2}}{(1 - \alpha_{t})\beta_t}\E\bracket{\nabla V_{t+1}^{*}\parenthesis{\frac{1}{\sqrt{\alpha_t}}\parenthesis{y_1 + (1 - \alpha_t)u_t^*(y_1)} + \sigma_t W_t)}} \\ & \quad - \frac{\sqrt{\alpha_{t}}\sigma_{t}^{2}}{(1 - \alpha_{t})\beta_t}\E\bracket{\nabla V_{t+1}^{*}\parenthesis{\frac{1}{\sqrt{\alpha_t}}\parenthesis{y_2 + (1 - \alpha_t)u_t^*(y_2)} + \sigma_t W_t)}}.
\end{align*}
Utilizing the Lipschitz condition of $s_t^{\rm pre}$ and $\nabla V_{t+1}^{*}$, we obtain
\begin{align*}
    \norm{u_t^*(y_1) - u_t^*(y_2)}_2 & \leq \norm{s_{t}^{\rm pre}(y_1) - s_{t}^{\rm pre}(y_2) }_2  + \frac{\sqrt{\alpha_{t}}\sigma_{t}^{2}}{(1 - \alpha_{t})\beta_t}\E\Big[\Big\|\nabla V_{t+1}^{*}\parenthesis{\frac{1}{\sqrt{\alpha_t}}\parenthesis{y_1 + (1 - \alpha_t)u_t^*(y_1)} + \sigma_t W_t)} \\ & \qquad\qquad\qquad - \nabla V_{t+1}^{*}\parenthesis{\frac{1}{\sqrt{\alpha_t}}\parenthesis{y_2 + (1 - \alpha_t)u_t^*(y_2)} + \sigma_t W_t)}\Big\|_2\Big] \\ & \leq L_{0, t}^s\norm{y_1 - y_2}_2 \\ & \quad + \frac{\sqrt{\alpha_{t}}\sigma_{t}^{2}}{(1 - \alpha_{t})\beta_t}L_{1, t+1}^{V^*}\norm{\frac{1}{\sqrt{\alpha_t}}\parenthesis{y_1 + (1 - \alpha_t)u_t^*(y_1)} - \frac{1}{\sqrt{\alpha_t}}\parenthesis{y_2 + (1 - \alpha_t)u_t^*(y_2)}}_2 \\ & \leq L_{0, t}^s\norm{y_1 - y_2}_2  + \frac{\sqrt{\alpha_{t}}\sigma_{t}^{2}}{(1 - \alpha_{t})\beta_t}L_{1, t+1}^{V^*}\parenthesis{\frac{1}{\sqrt{\alpha_t}}\norm{y_1 - y_2}_2 + \frac{1 - \alpha_t}{\sqrt{\alpha_t}}\norm{u_t^*(y_1) - u_t^*(y_2)}_2}.
\end{align*}
Equivalently, we have 
\begin{align*}
      \parenthesis{1 - \frac{\sigma_t^2}{\beta_t}L^{V^*}_{1, t+1}}\norm{u_t^*(y_1) - u_t^*(y_2)}_2  & \leq \parenthesis{L_{0, t}^s + \frac{\sigma_t^2 L^{V^*}_{1, t+1}}{(1- \alpha_t)\beta_t}}\norm{y_1 - y_2}_2.
\end{align*}
Since $1 - \frac{\sigma_t^2}{\beta_t}L^{V^*}_{1, t+1}\geq \lambda_t > 0$, we established the Lipschitz condition of the optimal control $u_t^*$:
\begin{align*}
    \norm{u_t^*(y_1) - u_t^*(y_2)}_2 \leq \lambda_{t}^{-1}\parenthesis{L_{0, t}^s + \frac{1 - \lambda_t}{1- \alpha_t}}\norm{y_1 - y_2}_2 = L_{0, t}^{u^*}\norm{y_1 - y_2}_2.,
\end{align*}
where we recall $L_{0, t}^{u^*}$ defined in \eqref{eq:Lu0star}.

%
%
%
%

%
%
%
%
%
%
%
%
%
%

\underline{Step 2: Differentiability of $ u_{t}^{*} $.} We now argue that $ u_t^*(y)$ is differentiable for all $y\in \mathbb{R}^d$. Let $h \in \R^d$ be an arbitrary non-zero vector and let $y \in \R^d$ be fixed. Since $s_t^{\rm pre}$ is differentiable, we have
\begin{align}
    s_t^{\rm pre}(y + h) - s_t^{\rm pre}(y) = \nabla s_t^{\rm pre}(y) h + o(\norm{h}_2). \label{eq:s-diff}
\end{align}
Moreover, by the inductive hypothesis, $\nabla V_{t+1}^*(y)$ is Lipschitz and thus $\nabla^2 V_{t+1}^*(y)$ exists for almost all $y$ and is bounded. \cite[Theorem 2.27]{folland1999real} implies the mapping $z \mapsto \E\bracket{\nabla V^{*}_{t+1}(z + \sigma_t W_t)}$ is differentiable everywhere and its derivative is given by
\begin{align}\label{eq:derivative nabla V}
    \frac{\partial}{\partial z}\E\bracket{\nabla V^{*}_{t+1}(z + \sigma_t W_t)} = \E\bracket{\nabla^2 V^{*}_{t+1}(z + \sigma_t W_t)}. 
\end{align}
Let $z = \frac{1}{\sqrt{\alpha_t}}\parenthesis{y + (1 - \alpha_t)u_t^*(y)} $ and $k = \frac{1}{\sqrt{\alpha_t}}(h + (1 - \alpha_t)(u_t^{*}(y+h) - u_t^{*}(y))$. Eq.~\eqref{eq:derivative nabla V} implies
\begin{align}
    \E\bracket{\nabla V^{*}_{t+1}(z + k + \sigma_t W_t)} - \E\bracket{\nabla V^{*}_{t+1}(z + \sigma_t W_t)} = \E\bracket{\nabla^2 V_{t+1}^*(z + \sigma_t W_t)}k + o(\norm{k}_2). \label{eq:nabla-v*-diff}
\end{align}
Since $u_t^*$ is Lipschitz, we have $k = \mcal{O}(\norm{h}_2)$ as $\norm{h}_2 \to 0$. Combining \eqref{eq:s-diff} and \eqref{eq:nabla-v*-diff}, we obtain
\begin{align*}
    u_t^*(y+h) - u_t^*(y) & = s_{t}^{\rm pre}(y+h) - s_{t}^{\rm pre}(y)  + \frac{\sqrt{\alpha_{t}}\sigma_{t}^{2}}{(1 - \alpha_{t})\beta_t}\E\bracket{\nabla V_{t+1}^{*}\parenthesis{\frac{1}{\sqrt{\alpha_t}}\parenthesis{(y+h) + (1 - \alpha_t)u_t^*(y+h)} + \sigma_t W_t)}} \\ & \quad - \frac{\sqrt{\alpha_{t}}\sigma_{t}^{2}}{(1 - \alpha_{t})\beta_t}\E\bracket{\nabla V_{t+1}^{*}\parenthesis{\frac{1}{\sqrt{\alpha_t}}\parenthesis{y + (1 - \alpha_t)u_t^*(y)} + \sigma_t W_t)}} \\ & = \nabla s_t^{\rm pre}(y) h  + \frac{\sqrt{\alpha_{t}}\sigma_{t}^{2}}{(1 - \alpha_{t})\beta_t}\mcal{H}_{t+1}(y)\parenthesis{\frac{1}{\sqrt{\alpha_t}}(h + (1 - \alpha_t)(u_t^{*}(y+h) - u_t^{*}(y))} + o(\norm{h}_2),
\end{align*}
where $\mcal{H}_{t+1}(y) \coloneqq \E\bracket{\nabla^2 V_{t+1}^{*}\left(\frac{1}{\sqrt{\alpha_t}}\parenthesis{y + (1 - \alpha_t)u_t^*(y)} + \sigma_t W_t\right)}  $.
Re-arranging the terms leads to
\begin{align*}
    u_t^*(y+h) - u_t^*(y) = \parenthesis{I_d - \frac{\sigma_t^2}{\beta_t}\mcal{H}_{t+1}(y)}^{-1}\parenthesis{\nabla s_t^{\rm pre}(y) + \frac{\sigma_t^2}{(1 - \alpha_t)\beta_t}\mcal{H}_{t+1}(y)}h + o(\norm{h}_2),
\end{align*}
which proves that $u_t^*(y)$ is differentiable for any $y\in \R^d$ and its derivative is given by 
\begin{align}\label{eq:grad u*-re}
    \nabla u_t^*(y) = \parenthesis{I_d - \frac{\sigma_t^2}{\beta_t}\mcal{H}_{t+1}(y)}^{-1}\parenthesis{\nabla s_t^{\rm pre}(y) + \frac{\sigma_t^2}{(1 - \alpha_t)\beta_t}\mcal{H}_{t+1}(y)}.
\end{align}

\underline{Step 3: Lipschitz condition of $  \nabla u_{t}^{*}$.} Next, we show that $\nabla u_t^*$ is Lipschitz. Taking the derivative of both sides of~\eqref{eq:u*(y)-re},
\begin{align}
    \nabla u_{t}^{*}(y) = \nabla s_{t}^{\rm pre}(y) + \frac{\sqrt{\alpha_{t}}\sigma_{t}^{2}}{(1 - \alpha_{t})\beta_t}\cdot\frac{\partial}{\partial y} \E\bracket{\nabla {V}^*_{t+1}\left(\frac{1}{\sqrt{\alpha_t}}\parenthesis{y + (1 - \alpha_t)u_t^{*}(y)} + \sigma_t W_t\right)}. \label{eq:u-diff}
\end{align}
For any $ z \in \R^{d} $, the integration by parts formula implies that
\begin{align}\label{eq:magic-1}
    \frac{\partial}{\partial z}\E\bracket{\nabla {V}^*_{t+1}\left(z + \sigma_t W_t\right)} = \E\bracket{\nabla^{2} {V}^*_{t+1}\left(z + \sigma_t W_t\right)} = \E\bracket{\nabla {V}^*_{t+1}\left(z + \sigma_t W_t\right)\frac{W_t\top}{\sigma_t}}.
\end{align}
Substituting $ z = \frac{1}{\sqrt{\alpha_t}}\parenthesis{y + (1 - \alpha_t)u_t^{*}(y)} $ and applying the chain rule, we obtain
\begin{align*}
     \frac{\partial}{\partial y}\E\bracket{\nabla {V}^*_{t+1}\left(\frac{1}{\sqrt{\alpha_t}}\parenthesis{y + (1 - \alpha_t)u_t^{*}(y)} + \sigma_t W_t\right)}  = \mcal{W}_{t+1}(y)\mcal{U}_t(y),
\end{align*}
where $\mcal{W}_{t+1}(y)  \coloneqq  \E\bracket{\nabla {V}^*_{t+1}\left(\frac{1}{\sqrt{\alpha_t}}\parenthesis{y + (1 - \alpha_t)u_t^{*}(y)} + \sigma_t W_t\right)\frac{W_t^\top}{\sigma_t}}$ and $\mathcal{U}_t(y) \coloneqq \frac{1}{\sqrt{\alpha_t}}\parenthesis{I_d + (1 - \alpha_t) \nabla u_t^*(y)}$. It follows from \eqref{eq:u-diff} that
\begin{align}
    \nabla u_{t}^{*}(y)  & = \nabla s_{t}^{\rm pre}(y)  + \frac{\sqrt{\alpha_{t}}\sigma_{t}^{2}}{(1 - \alpha_{t})\beta_t} \mcal{W}_{t+1}(y)\mcal{U}_t(y). \label{eq:u-diff-1}
\end{align}
For any $y_1$ and $y_2$ in $\R^d$, Eq.~\eqref{eq:u-diff-1} implies that %
\begin{align}
\nabla u_t^{*} (y_1) - \nabla u_t^{*} (y_2)  & = \nabla s_{t}^{\rm pre}(y_{1}) - \nabla s_{t}^{\rm pre}(y_{2})  +  \frac{\sqrt{\alpha_{t}}\sigma_{t}^{2}}{(1 - \alpha_{t})\beta_t} \big(\mcal{W}_{t+1}(y_1)\mcal{U}_t(y_1)  - \mcal{W}_{t+1}(y_2)\mcal{U}_t(y_2)\big). \label{eq:u-lip-1}
\end{align}
Note that for any $ y \in \R^d$ we have
\begin{align}
    & \norm{\mcal{W}_{t+1}(y)}_{2} = \norm{\mcal{H}_{t+1}(y)}_2 \leq L_{1, t+1}^{V^*}, \label{eq:upper-bound-magic-1} \\ 
    & \norm{\mcal{U}_t(y)}_2 \leq \frac{1}{\sqrt{\alpha_{t}}}\parenthesis{1 + (1 - \alpha_{t})L_{0,t}^{u^{*}}}, \label{eq:upper-bound-nabla-u-n-1}
\end{align}
where \eqref{eq:upper-bound-magic-1} holds by applying the identity \eqref{eq:magic-1}, and \eqref{eq:upper-bound-nabla-u-n-1} is a consequence of the Lipschitz condition of $ u_{t}^{*} $.
Moreover, for any $ y_{1} $ and $ y_{2} $ in $\R^d$, the Lipschitz conditions of $ \nabla {V}^*_{t+1} $ and~$ u_{t}^{*} $ imply that
\begin{align}
     \norm{\mcal{W}_{t+1}(y_1) - \mcal{W}_{t+1}(y_2)}_2 \nonumber &  \leq \frac{\E\bracket{\norm{W_{t}}_2}}{\sigma_{t}}L_{1, t+1}^{{V}^*}\frac{1}{\sqrt{\alpha_{t}}}\parenthesis{\norm{y_{1} - y_{2}}_2 + (1 - \alpha_{t})L_{0, t}^{u^{*}}\norm{y_{1} - y_{2}}_2} \nonumber \\ & = \frac{\E\bracket{\norm{W_{t}}_2}}{\sqrt{\alpha_{t}}\sigma_{t}}L_{1, t+1}^{V^{*}}\parenthesis{1 + (1 - \alpha_{t})L_{0, t}^{u^{*}}}\norm{y_{1} - y_{2}}_2. \label{eq:Lip-nabla-V-w-1}
\end{align}
Furthermore, the definition of $\mcal{U}_t$ implies
\begin{align}
    \norm{\mcal{U}_t(y_1) - \mcal{U}_t(y_2)}_2 & = \frac{1 - \alpha_{t}}{\sqrt{\alpha_{t}}}\norm{\nabla u_{t}^{*}(y_{1}) - \nabla u_{t}^{*}(y_{2})}_2 \label{eq:Lip-nabla-u-n-1}.
\end{align}
Combining \eqref{eq:u-lip-1}--\eqref{eq:Lip-nabla-u-n-1} together, we have
\begin{align*}
    \norm{\nabla u_{t}^{*}(y_{1}) - \nabla u_{t}^{*}(y_{2})}_2 & \leq \norm{\nabla s_{t}^{\rm pre}(y_{1}) - \nabla s_{t}^{\rm pre}(y_{2})}_2  +  \frac{\sqrt{\alpha_{t}}\sigma_{t}^{2}}{(1 - \alpha_{t})\beta_t} \norm{\mcal{W}_{t+1}(y_1)  - \mcal{W}_{t+1}(y_2)}_2\norm{\mcal{U}_t(y_1)}_2 \\ & \qquad + \frac{\sqrt{\alpha_{t}}\sigma_{t}^{2}}{(1 - \alpha_{t})\beta_t} \norm{\mcal{W}_{t+1}(y_2)}_2\norm{\mcal{U}_t(y_1) - \mcal{U}_t(y_2)}_2 
    \\ & \leq L_{1, t}^{s}\norm{y_{1} - y_{2}}_2 + \frac{\sqrt{\alpha_{t}}\sigma_{t}^{2}}{(1 - \alpha_{t})\beta_t}\bigg(L_{1, t+1}^{V^{*}}\frac{1 - \alpha_{t}}{\sqrt{\alpha_{t}}} \norm{\nabla u_{t}^{*}(y_{1}) - \nabla u_{t}^{*}(y_{2})}_2 \\ & \qquad + \frac{1}{\alpha_{t}}\parenthesis{1 + (1 - \alpha_{t})L_{0,t}^{u^{*}}}^{2}\frac{\E\bracket{\norm{W_{t}}_2}}{\sigma_{t}}L_{1, t+1}^{V^{*}}\norm{y_{1} - y_{2}}_2\bigg).
\end{align*}
Since $ \frac{\sigma_{t}^{2}}{\beta_{t}}L_{1, t+1}^{V^{*}} \leq 1 - \lambda_{t} $, we deduce that
\begin{align*}
    \norm{\nabla u_{t}^{*}(y_{1}) - \nabla u_{t}^{*}(y_{2})}_2 & \leq \lambda_{t}^{-1}\parenthesis{L_{1, t}^{s} +  \frac{\E\bracket{\norm{W_{t}}_2}(1 - \lambda_t)}{(1 - \alpha_t)\sqrt{\alpha_{t}}\sigma_{t}}\parenthesis{1 + (1 - \alpha_{t})L_{0,t}^{u^{*}}}^{2}}\norm{y_1 - y_2}_2 \\ &  = L_{1, t}^{u^*}\norm{y_1 - y_2}_2,
\end{align*}
where we recall $L_{1, t}^{u^*}$ defined in \eqref{eq:Lu1star}.


\underline{Step 4: Lipschitz conditions of $ V_{t}^{*} $ and $ \nabla V_{t}^{*}$.} Finally, we turn to prove the Lipschitz and gradient Lipschitz conditions of $V_t^*$. Plugging $u_t^*$ into the Bellman equation \eqref{eq:bellman-V*}, we have
\begin{align}\label{eq:bellman-u*}
 V_t^*(y) & = \E\bracket{V^{*}_{t+1}\left(\frac{1}{\sqrt{\alpha_t}}\parenthesis{y + (1 - \alpha_t)u_t^*(y)} + \sigma_t W_t\right)} - \beta_t \frac{(1 - \alpha_{t})^{2}}{2\alpha_{t}\sigma_{t}^{2}}\norm{u_t^*(y) - s^{\rm pre}_{t}(y)}_{2}^{2}.
 %
\end{align}
Since $V_{t+1}^*$ and $s_t^{\rm pre}$ are differentiable with Lipschitz gradients and $u_t^*$ is differentiable, we know that $V_{t}^{*}$ is differentiable and 
\begin{align*}
 \nabla V_t^*(y) &= \frac{\partial}{\partial y} \E\Bigg[V^{*}_{t+1}\left(\frac{1}{\sqrt{\alpha_t}}\parenthesis{y + (1 - \alpha_t)u_t^*(y)} + \sigma_t W_t\right)- \beta_t \frac{(1 - \alpha_{t})^{2}}{2\alpha_{t}\sigma_{t}^{2}}\norm{u_t^*(y) - s^{\rm pre}_{t}(y)}_{2}^{2}\Bigg]  \\
 &= \frac{1}{\sqrt{\alpha_t}}\parenthesis{I_d 
 + (1 - \alpha_t) \nabla u_t^*(y)}^\top\E\bracket{\nabla V_{t+1}^{*}\left(\frac{1}{\sqrt{\alpha_t}}\parenthesis{y + (1 - \alpha_t)u_t^*(y)} + \sigma_t W_t\right)}\\
 & \quad - \beta_t \frac{(1 - \alpha_{t})^{2}}{\alpha_{t}\sigma_{t}^{2}}(\nabla u_t^*(y) - \nabla s^{\rm pre}_{t}(y))^\top(u_t^* - s^{\rm pre}_{t}(y)).
\end{align*}
Define $\mathcal{G}_{t+1}(y) \coloneqq \E\bracket{\nabla V_{t+1}^{*}\left(\frac{1}{\sqrt{\alpha_t}}\parenthesis{y + (1 - \alpha_t)u_t^*(y)} + \sigma_t W_t\right)} $. It follows that 
\begin{eqnarray} 
 \nabla V_t^*(y) 
 &=&\frac{1}{\sqrt{\alpha_t}}\parenthesis{I_d 
 + (1 - \alpha_t) \nabla u_t^*(y)}^\top\mathcal{G}_{t+1}(y)- \beta_t \frac{(1 - \alpha_{t})^{2}}{\alpha_{t}\sigma_{t}^{2}}(\nabla u_t^*(y) - \nabla s^{\rm pre}_{t}(y))^\top\left(  \frac{\sqrt{\alpha_{t}}\sigma_{t}^{2}}{(1 - \alpha_{t})\beta_t}\mathcal{G}_{t+1}(y)\right)\nonumber\\
 &=&\frac{1}{\sqrt{\alpha_t}}\parenthesis{I_d 
 + (1 - \alpha_t) \nabla s_t^{\rm pre}(y)}^\top\mathcal{G}_{t+1}(y) \label{eq:grad-V*},
\end{eqnarray}
where \eqref{eq:u*(y)-re} is applied to obtain the first equality. Next, we use \eqref{eq:grad-V*} to establish the Lipschitz conditions of $ V_{t}^{*} $ and $\nabla V_t^{*}$. The Lipschitz condition of $\nabla V_{t+1}^*$ implies $\norm{\mcal{G}_{t+1}(y)}_2 \leq L_{0, t+1}^{V^*}$ for all $y \in \R^d$. Based on \eqref{eq:grad-V*}, we have
\begin{align*}
\norm{\nabla V_{t}^{*}(y)}_{2} \leq \frac{1}{\sqrt{\alpha_t}}(1 + (1 - \alpha_t)L_{0, t}^{s})
L_{0, t+1}^{V^*} = L_{0, t}^{V^*},
\end{align*}
which proves the Lipschitz condition of $ V_t^* $. Next, we prove the gradient Lipschitz condition of $ V_t^* $. For ease of exposition, we denote $\mathcal{S}_t(y) \coloneqq \frac{1}{\sqrt{\alpha_t}}\parenthesis{I_d 
+ (1 - \alpha_t) \nabla s_t^{\rm pre}(y)}$. We now use \eqref{eq:grad-V*} to show the Lipschitz condition of $\nabla V_t^*$. We first note that $ \nabla \mcal{G}_{t+1}(y) = \mcal{H}_{t+1}(y)\mcal{U}_t(y) $ is well-defined at every $y \in \R^d$, and thus
\begin{align*}
\norm{\nabla \mcal{G}_{t+1}(y)}_{2} \leq \norm{\mcal{H}_{t+1}(y)}_{2}\norm{\mcal{U}_t(y)}_{2} \leq \frac{1}{\sqrt{\alpha_{t}}}(1 + (1 - \alpha_{t})L_{0, t}^{u^*})L_{1, t+1}^{V^*}.
\end{align*}
 With the Lipschitz condition of $ \mcal{G}_{t+1} $ in hand,  for any $y_1$ and $y_2$ in $\R^d$, we have that
\begin{align*}
\norm{\nabla V_{t}^{*}(y_{1}) - \nabla V_{t}^{*}(y_{2})}_{2} & \leq \norm{\mcal{S}_t(y_{1})^{\top}\mcal{G}_{t+1}(y_{1}) - \mcal{S}_t(y_{2})^{\top}\mcal{G}_{t+1}(y_{2})}_2 \\ & \leq \norm{\mcal{S}_t(y_{1})^{\top}\parenthesis{\mcal{G}_{t+1}(y_{1}) - \mcal{G}_{t+1}(y_{2})}}_2 + \norm{\parenthesis{\mcal{S}_t(y_{1}) - \mcal{S}_t(y_{2})}^{\top}\mcal{G}_{t+1}(y_{2})}_{2} \\ & \leq \norm{\mcal{S}_t(y_{1})}_{2}\norm{\mcal{G}_{t+1}(y_{1}) - \mcal{G}_{t+1}(y_{2})}_{2} + \norm{\mcal{S}_t(y_{1}) - \mcal{S}_t(y_{2})}_{2}\norm{\mcal{G}_{t+1}(y_{2})}_{2} \\ & \leq \frac{1}{\alpha_{t}}(1 + (1 - \alpha_{t})L_{0, t}^s)(1 + (1 - \alpha_{t})L_{0, t}^{u^*})L_{1,t+1}^{V^*}\norm{y_{1} - y_{2}}_{2} \\ & \quad + \frac{1 - \alpha_{t}}{\sqrt{\alpha_{t}}}L_{1,t}^s L_{0, t+1}^{V^*}\norm{y_{1} - y_{2}}_{2},
\end{align*}
where the last equation uses the Lipschitz conditions of $s_t^{\rm pre}$, $\nabla s_t^{\rm pre}$, $V_{t+1}^*$ and $\nabla V_{t+1}^*$. Consequently, we have
\begin{align*}
    \norm{\nabla V_{t}^{*}(y_{1}) - \nabla V_{t}^{*}(y_{2})}_{2} \leq L_{1, t}^{V^*}\norm{y_1 - y_2}_2,
\end{align*}
where we recall $L_{1, t}^{V^*}$ defined in \eqref{eq:LV1star}. In other words, $L_{1, t}^{V^*}$ is indeed the Lipschitz constant of $\nabla V_t^*$. This completes the proof. 


\end{proof}

\begin{thebibliography}{99}
\bibitem[8]{black2023training} Kevin Black, Michael Janner, Yilun Du, Ilya Kostrikov, and Sergey Levine. Training diffusion models with reinforcement learning. In \emph{The Twelfth International Conference on Learning Representations}, 2024. URL \url{https://openreview.net/forum?id=YCWjhGrJFD}.
\bibitem[12]{dhariwal2021diffusion} Prafulla Dhariwal and Alexander Nichol. {Diffusion models beat GANs on image synthesis}. \emph{Advances in neural information processing systems}, 34:\penalty0 8780--8794, 2021.
\bibitem[18]{fan2024reinforcement} Ying Fan, Olivia Watkins, Yuqing Du, Hao Liu, Moonkyung Ryu, Craig Boutilier, Pieter Abbeel, Mohammad Ghavamzadeh, Kangwook Lee, and Kimin Lee. Reinforcement learning for fine-tuning text-to-image diffusion models. \emph{Advances in Neural Information Processing Systems}, 36, 2024.
\bibitem[22]{folland1999real} Gerald~B Folland. \emph{Real analysis: modern techniques and their applications}, volume~40. John Wiley \& Sons, 1999.
\bibitem[26]{gao2024reward} Xuefeng Gao, Jiale Zha, and Xun~Yu Zhou. Reward-directed score-based diffusion models via q-learning. \emph{arXiv preprint arXiv:2409.04832}, 2024.
\bibitem[30]{han2024neural} Yinbin Han, Meisam Razaviyayn, and Renyuan Xu. Neural network-based score estimation in diffusion models: Optimization and generalization. In \emph{International Conference on Learning Representations}, 2024. URL \url{https://openreview.net/forum?id=h8GeqOxtd4}.
\bibitem[32]{Ho20} Jonathan Ho, Ajay Jain, and Pieter Abbeel. Denoising diffusion probabilistic models. In \emph{Advances in Neural Information Processing Systems}, volume~33, pages 6840--6851, 2020.
\bibitem[34]{hyvarinen2005estimation} Aapo Hyv{\"a}rinen and Peter Dayan. Estimation of non-normalized statistical models by score matching. \emph{Journal of Machine Learning Research}, 6\penalty0 (4), 2005.
\bibitem[38]{li2023towards} Gen Li, Yuting Wei, Yuxin Chen, and Yuejie Chi. Towards non-asymptotic convergence for diffusion-based generative models. In \emph{The Twelfth International Conference on Learning Representations}, 2024. URL \url{https://openreview.net/forum?id=4VGEeER6W9}.
\bibitem[49]{ramesh2022hierarchical} Aditya Ramesh, Prafulla Dhariwal, Alex Nichol, Casey Chu, and Mark Chen. Hierarchical text-conditional image generation with clip latents. \emph{arXiv preprint arXiv:2204.06125}, 1\penalty0 (2):\penalty0 3, 2022.
\bibitem[50]{rombach2022high} Robin Rombach, Andreas Blattmann, Dominik Lorenz, Patrick Esser, and Bj{\"o}rn Ommer. High-resolution image synthesis with latent diffusion models. In \emph{Proceedings of the IEEE/CVF conference on computer vision and pattern recognition}, pages 10684--10695, 2022.
\bibitem[55]{song2020score} Yang Song, Jascha Sohl-Dickstein, Diederik~P Kingma, Abhishek Kumar, Stefano Ermon, and Ben Poole. Score-based generative modeling through stochastic differential equations. In \emph{International Conference on Learning Representations}, 2021. URL \url{https://openreview.net/forum?id=PxTIG12RRHS}.
\bibitem[66]{watson2023novo} Joseph~L Watson, David Juergens, Nathaniel~R Bennett, Brian~L Trippe, Jason Yim, Helen~E Eisenach, Woody Ahern, Andrew~J Borst, Robert~J Ragotte, Lukas~F Milles, et~al. De novo design of protein structure and function with rfdiffusion. \emph{Nature}, 620\penalty0 (7976):\penalty0 1089--1100, 2023.
\bibitem[71]{zhang2001some} Jianfeng Zhang. \emph{Some fine properties of backward stochastic differential equations}. Purdue University, 2001.
\end{thebibliography}
\end{document}
